% !TeX root = ../../main.tex

Dieser Abschnitt handelt von der Fehlerbehandlung in Datenbanken.
Dafür wird kurz beschrieben, wie die Fehler bisher erkannt und verarbeitet werden.

Bisher behandelt eine Datenbank Fehler nicht von selbst.
Sobald ein Fehler erkannt wird, wird eine Transaktion abgebrochen und er Client bekommt eine Fehlermeldung.
Dieser muss dann entscheiden wie mit diesem Fehler umgegangen wird.
Bei einigen Fehlern ist dies sinnvoll, bei anderen ist es erhöhter Aufwand für den Client, beziehungsweise für
den Nutzer.

Es wurde sich für Recovery die Mühe gegeben, bestimmte Mechanismen einzuführen, um die Wiederherstellung nach
einem Absturz zu gewährleisten.
So gibt es das in~\ref{sec:postgresql} erwähnte WAL, welches auf \texttt{UNDO} und \texttt{REDO} basiert.
Für \texttt{UNDO} und \texttt{REDO} gibt es aber noch unterteilungen.
Es gibt \texttt{Transaction UNDO}, welches alle Effekte der abgebrochenen Transaktion entfernt und
das \texttt{Global UNDO}, welches alle Effekte von allen abgebrochenen Transaktionen rückgängig macht.
Das \texttt{Global UNDO} wird häufig nach einem Systemabsturz verwendet.
\texttt{Partial REDO} wird ausgeführt, wenn Ergebnisse von Transaktionen noch nicht in der Datenbank standen,
als das System abgestürzt ist.
Dadurch werden die Transaktionen wiederholt und die Ergebnisse werden in die Datenbank geschrieben.
Bei \texttt{Global REDO} handelt es sich um einen Sonderfall, der in der Regel nur bei \emph{Media Failure}\footnote{Ein
	Media Failure tritt auf, wenn Speicher verloren oder kaputt geht.} auftritt.
Da nach Media Failure eine ältere Version der Datenbank vorliegt, werden durch \texttt{Global REDO} alle Transaktionen
bis zu dem Zeitpunkt des Failures nachgeholt, damit diese wieder in der Datenbank stehen.
Des Weiteren werden Checkpoints verwendet, um den Recovery Aufwand zu optimieren und limitieren~\cite{Haerder1983}.

Weitere interne Fehlerbehandlung passiert durch \emph{Optimistic Concurrency Control (OCC)},
\emph{Pessimistic Concurrency Control (PCC)} und \emph{Multi-Version Concurrency Control (MVCC)}\@.
OCC funktioniert, in dem es während einer Transaktion die Locks in einer kurzen Zeit bekommt und wieder freigibt.
Die Locks werden direkt vor einer Leseoperation erworben und sobald diese fertig ist, wieder zugänglich gemacht.
Für eine Schreiboperation werden die Lock auch erst direkt vor der Operation erworben, aber bis zum Ende der
Transaktion gehalten.
Bei PCC werden die Locks am Anfang der Transaktion erworben und erst beim Abschluss wieder freigegeben, dadurch kann
eine Ressource lange gesperrt werden~\cite{IBMConcurrencyControl}.
Das von PostgreSQL verwendete MVCC zeigt jeder Transaktion einen Snapshot der Datenbank.
Dies sorgt dafür, dass die Locks zum Lesen nicht mit den Locks zum Schreiben keine Konflikte verursachen~\cite{PostgreSQLMVCC}.